\begin{corollary}[5.4.2]
\label{IV.2.s5.4.2-fr}
Si $X$, $Y$ son dos preesquemas localmente noetherianos y $f:X\to Y$ es un
morfismo finito \textup{(por tanto} cerrado\textup{)}, se tiene
$\dim(X)=\dim(f(X))$. Si además $f$ es sobreyectivo, se tiene
$\dim(X)=\dim(Y)$.
\end{corollary}

\begin{remarks}[5.4.3]
\label{IV.2.s5.4.3-fr}
\begin{enumerate}
\item[(i)] Se ha visto en (4.1.2) que, si $X$ e $Y$ son preesquemas
localmente de tipo finito sobre un cuerpo $k$ y $f$ es un $k$-morfismo, la
desigualdad $\dim(Y)\leq\dim(X)$ se verifica ya cuando $f$ es \emph{cuasicompacto y
dominante}. Por el contrario, se puede tener $\dim(Y)>\dim(X)$ incluso cuando
$f$ es \emph{de tipo finito}, \emph{biyectivo} y una \emph{inmersión local}, si se supone solamente
que $X$ e $Y$ son localmente noetherianos. Por ejemplo, sean $A$ un anillo de
valoración discreta, $K$ su cuerpo de fracciones y $k$ su cuerpo residual; si
$Y=\operatorname{Spec}(A)$, y si $a$ es el punto genérico de $Y$ y $b$ su
punto cerrado, se tiene, pues, $\boldsymbol{k}(a)=K$, $\boldsymbol{k}(b)=k$. Sea entonces $X$ el
preesquema \emph{suma} de $\operatorname{Spec}(K)$ y de $\operatorname{Spec}(k)$, y
$f:X\to Y$ el morfismo que coincide en $\operatorname{Spec}(K)$ y
$\operatorname{Spec}(k)$ con los respectivos morfismos canónicos (\textbf I, 2.4.5);
es claro que $f$ es una inmersión local biyectiva, pues $\{a\}$ es abierto en
$Y$; por otra parte, $f$ es de tipo finito, ya que $K=A[\pi^{-1}]$, donde
$\pi$ es un uniformizante de $A$, y por tanto $K$ es una $A$-álgebra de tipo
finito. Sin embargo, se tiene $\dim(X)=0$ y $\dim(Y)=1$.

\item[(ii)] Si $A$ y $B$ son dos anillos noetherianos tales que $A\subset B$
y $B$ es una $A$-álgebra \emph{finita}, el corolario (5.4.2) muestra de nuevo que
$\dim(B)=\dim(A)$ (\textbf 0, 16.1.5). Supongamos además que $A$ es un anillo \emph{local}
noetheriano; entonces $B$ es un anillo noetheriano \emph{semi-local} (Bourbaki,
\emph{Alg. comm.}, cap. IV, \S~2, n\textsuperscript{o}~5, cor.~3 de la
prop.~9); si $\mathfrak n_i$ $(1\leq i\leq r)$ son los ideales maximales de
$B$, se tiene, pues,
\begin{equation}
\dim(A)=\sup_i\dim(B_{\mathfrak n_i}).
\tag{5.4.3.1}\label{IV.2.s5.4.3.1-fr}
\end{equation}
Se observará que, en estas condiciones, no se tiene necesariamente
$\dim(B_{\mathfrak n_i})=\dim(A)$ para \emph{todo} $i$: para verlo basta reemplazar
$B$ por el producto directo de $B$ y del cuerpo residual $k$ de $A$. Pero
existen incluso ejemplos en los que $A$ y $B$ son además \emph{anillos íntegros de
dimensión} $2$ y en los que algunos de los $B_{\mathfrak n_i}$ tienen dimensión
$<\dim(A)$ (5.6.11). Veremos, sin embargo, más adelante ((5.6.4) y (5.6.10))
que este último fenómeno no puede presentarse cuando se supone que $A$ es
cociente de un anillo local regular.
\end{enumerate}
\end{remarks}

\subsection{Fórmula de las dimensiones para un morfismo de tipo finito.}
\label{IV.2.s5.5-fr}
